% Source: https://artofproblemsolving.com/community/c3h1351874_math_league_question_19
% Stable identity: aops-topic-1351874-q19
Since $\angle C=\angle D=90^\circ$ and $BC=CD=DE=4$, the quadrilateral $BCDE$ is a square of side $4$, so $BE=4$ and $\angle CBE=\angle DEB=90^\circ$. Then $\angle ABE=150^\circ-90^\circ=60^\circ$ and likewise $\angle AEB=60^\circ$, and $AB=AE=BE=4$: triangle $ABE$ is equilateral and lies outside the square. (The angle at $A$ is $540^\circ-2\cdot90^\circ-2\cdot150^\circ=60^\circ$, as it should be.)

\begin{asy}
usepackage("amsmath");
size(4.5cm);
defaultpen(fontsize(10pt));
path ra(pair P,pair u,pair v,real s=0.2){return P+s*unit(u)--P+s*unit(u)+s*unit(v)--P+s*unit(v);}
pair C=(0,0),D=(4,0),Ee=(4,4),B=(0,4),A=(2,4+2*sqrt(3));
draw(A--B--C--D--Ee--cycle,linewidth(0.9));
draw(B--Ee,red+dashed);
draw(ra(C,D-C,B-C,0.4)); draw(ra(D,Ee-D,C-D,0.4));
label("$A$",A,N); label("$B$",B,W); label("$C$",C,SW); label("$D$",D,SE); label("$E$",Ee,E);
label("$4$",(C+D)/2,S); label("$4$",(B+C)/2,W); label("$4$",(D+Ee)/2,E);
label("$4$",(A+B)/2,NW); label("$4$",(A+Ee)/2,NE); label("$4$",(B+Ee)/2,S,red);
label("\mbox{}",(min(currentpicture,user=true)+max(currentpicture,user=true))/2,black); // keeps the text after the figure black when LaTeX typesets the labels inline
\end{asy}

So the area is
\[
[BCDE]+[ABE]=4^2+\dfrac{\sqrt3}{4}\cdot4^2=\boxed{16+4\sqrt3}.
\]
